\section{Introduction}
\label{sec:introduction}

Large orbital structures are becoming central to civil, scientific, and
commercial spaceflight. Space stations, assembled telescopes, logistics
platforms, and in-space manufacturing facilities all require periodic external
inspection to detect configuration changes, micrometeoroid damage, thermal
blanket degradation, docking-interface anomalies, and other faults that are
not easily observed from fixed onboard cameras. The corresponding planning
problem is simultaneously geometric and dynamic. A candidate view may provide
high visual coverage yet remain unsafe, expensive, or unreachable under
relative orbital motion.

Existing inspection planners often emphasize one side of this problem. Robotic
next-best-view methods provide mature tools for coverage, occlusion reasoning,
and sensor placement, but their motion cost is usually geometric. Relative
orbit-design methods provide dynamically meaningful paths, but they do not
always treat surface coverage, camera field of view, and target occlusion as
the planning objective. This separation is limiting for large structures.
Selecting views first and checking dynamics later can produce trajectories
that appear coverage-rich but violate clearance or require excessive input.

This paper addresses that gap by formulating orbital structure inspection as a
dynamics-aware coverage problem on a finite candidate-view graph. The planner
samples area-weighted surface targets, constructs camera-feasible views,
evaluates visibility under range, field-of-view, incidence-angle, and
line-of-sight constraints, and assigns each candidate transfer a
bounded-acceleration Clohessy--Wiltshire/Hill (CW/HCW) cost. View selection is
therefore driven by the coverage that a view adds and by the orbital effort,
clearance, and input saturation incurred to reach it.

The finite graph also gives a precise scope for optimality. For a fixed
candidate set, fixed visibility relation, and fixed feasible-transfer graph,
the inspection task can be written as a mixed-integer optimization problem. We
implement an exact dynamic-programming graph solver that returns a certificate
on reduced instances by minimizing CW dynamic cost subject to a required
coverage ratio and a view budget. The full ISS-mesh benchmark uses a scalable
dynamics-aware set-cover and CW-tour solver on the same graph. Thus, the paper
separates an implemented finite-graph optimality oracle from the scalable
planner used to generate the reported high-coverage trajectories.

The reported results show that dynamics awareness matters. On the ISS-mesh
benchmark, the proposed dynamics-aware tour achieves 98.33\% coverage with 21
selected views, \unit{21.592}{m/s} cumulative $\Delta v$, \unit{9.91}{m}
minimum clearance, and \unit{5.97}{s} planning time. Coverage-greedy reaches
the same coverage but requires input beyond the admissible bound and violates
clearance; nearest selection also reaches high coverage but violates
clearance. Safe-coverage and random-safe baselines remain feasible at high
coverage, but require substantially higher cumulative $\Delta v$, while the
fuel-focused baseline stops below the high-coverage target. The comparison
therefore exposes failure modes that are hidden by coverage metrics alone in
this deterministic base case. Robustness across target density, candidate
density, initial state, keep-out margin, and transfer duration is treated as a
separate validation requirement rather than inferred from one mesh instance.

ROS~2 Jazzy and Gazebo Harmonic are used as a replay and verification layer,
not as the source of orbital dynamics. The ROS stack publishes the chaser
state from the HCW dynamics model, records trajectory, attitude, coverage,
safety, and planner logs, and reconstructs the camera frustum used by the
offline planner. Gazebo provides visual context for the station and chaser.
This replay layer verifies that the planned visual geometry can be
instantiated in the integrated system while preserving the planner as the
scientific object of study.

The main contributions are:
\begin{enumerate}
  \item a finite candidate-graph formulation of orbital inspection that states
  the scope of the discrete optimization problem, including camera visibility,
  dwell coverage, CW/HCW transfer feasibility, clearance, and input limits;
  \item a two-mode solver for this graph: an exact dynamic-programming
  certificate on reduced finite graphs and a scalable dynamics-aware offline
  solver that couples set-cover gain with transfer cost and safety margins
  during view selection;
  \item a deterministic ISS-mesh base benchmark and validation matrix for
  target-density, candidate-density, initial-condition, safety-margin,
  transfer-duration, and ablation studies;
  \item a ROS~2/Gazebo replay layer that verifies camera frusta, logged
  attitudes, and trajectory geometry while keeping HCW dynamics as the source
  of quantitative evidence.
\end{enumerate}
